\section{Address counting and completion of the joint limit}\label{sec:address}

We complete the bad-grid estimate from Section~\ref{sec:endpoints}.  The
argument counts the addresses of an adaptively selected collection of bad
times and charges, on the original walk law, the necessary tests furnished
by the tangent passage exclusion.  The same estimate then supplies the
endpoint modulus needed to identify the whole distance graph.  A final
bounded calibration removes the deterministic metric factor along every
subsequence.

Throughout the counting argument we retain the exponent choices and the
lift of Section~\ref{sec:endpoints}.  In particular,
\[
 \frac14<\tau<\varrho<s<\sigma<\min(\xi,1),\qquad
 1<a<a_+<b_-<b,
 \qquad \mathcal I=[a_+N,b_-N].
\]
The interval from which the $d$ base times are sampled has width
$O(Te^{-N})$.  The moment order $d$ is fixed while the large parameters
tend to infinity.  We use the density bound
\eqref{eq:original-19} relative to the original disk law times counting
measure on the $d$ addresses.  The additional grid translations and probe
phases below are independent auxiliary randomness under both laws.

Fix an integer shell ratio $R$, put $g=\log R$, and write
\begin{equation}\label{addr:scales}
 r(j)=TR^{-j},\qquad
 P(j)=\frac{2}{\pi^2}r(j)\log r(j),\qquad
 0\le j\le J_0=\left\lfloor\frac{bN}{g}\right\rfloor.
\end{equation}
We may extend the range by a bounded number of levels.  All time scales
in use tend to infinity at least as a fixed positive power of the large
perimeter parameter.  Their logarithms are comparable, as established in
the endpoint reduction.  Consequently bounded shifts in $j$ have the
corresponding asymptotic exponential ratios.  Constants may depend on
the fixed exponent margins and, where indicated, on $d$ and $R$.

\subsection{Finite tests at sparse logarithmic levels}

Choose sparse probe indices $i$ with $gi\in\mathcal I$, separated by
$L$ in logarithmic units.  Here $L/g$ is an integer, $L\to\infty$
sufficiently slowly along the extraction, and $L=o(N)$.  Give the sparse
lattice a common uniformly sampled phase.  To a probe $i$ assign the
disjoint level interval
\begin{equation}\label{addr:assigned-levels}
                       i-L/g<j\le i-3.
\end{equation}
A fixed additional number of levels next to the probe may be omitted;
this number may depend on $d$ and $R$.

\begin{lemma}[Finite necessary tests]\label{addr:finite-tests}
With probability tending to one under the lift, all but $o(N/L)$ probe
levels admit, for each base, one of the following finite witnesses.
\begin{enumerate}[label=\textup{(\roman*)}]
\item A main singularity, or two main positive bits, with their
coordinates and time directions specified.
\item A main near-link: an integer state time $q\in[0,T]$ and a row
comparing one end height at the base with one end height at $q$, to
tolerance $e^{\sqrt L}P(i)$.  The time separation exceeds
$e^{2L}r(i)$.  On the interval between them, each indicated endpoint
coordinate, in the direction of the other endpoint, remains outward of
the comparison wall up to that tolerance.
\item One main positive bit, together with an auxiliary obstruction or
an auxiliary origin test.  The witness specifies its aperture rule,
seed and path indices, and its side and coordinate type.  An obstruction
requires a descending ladder record within $CP(i)$ of the appropriate
original-label depth in the auxiliary path's starting polygon.  Before
that path, the targeted arc has both margins greater than
$ce^{2L}P(i)$, its cyclic order persists, and both directed starting
distances to its middle label exceed $ce^{2L}P(i)$.  An origin test on an
accessible helper has the same long transmitted arc, starts within
$o(P(i))$ of its center, and requires a singularity or a positive bit
at the path origin.  For the count, this last positional bound may be
weakened to $O(P(i))$.
\end{enumerate}
At an assigned level $j$, a positive bit means that its coordinate stays
outward of its base or origin value, with tolerance $P(j+1)$, on the
slot in the specified direction between distances $r(j+1)$ and
$2r(j)$.  A singularity means that the maximal change within
$r(j+1)$ of the center is at least $P(j)$.  Bounded integer rounding
and fixed enlargements of these tolerances are allowed.  The coordinates
are $h,-b$, or the corresponding outward-growth coordinates of a helper.

If two bases at the same probe are more than $r(i)$ apart and make
auxiliary requests in the same filling $J$, their transmitted arc
positions can also be required to be separated on its circumference by
$ce^{2L}P(i)$.  After all positions have been specified, the remaining
tags have at most $\exp(O_d(1+N/L))$ possibilities.
\end{lemma}

\begin{proof}
Apply the typical-logarithmic-scale stationarity and absolute-continuity
alternatives of Section~\ref{sec:endpoints}, together with
Proposition~\ref{end:passage-exclusion}, first with all factors bounded.
On coupled further
extractions, the complement of the alternatives in the statement is a
null event.  The bounded factors may therefore increase sufficiently
slowly so that the resulting tests hold simultaneously on the shifted
levels assigned to a probe.  This is a diagonal choice of finite
witnesses.  It does not condition the tangent assertion on a prescribed
sparse phase: first sample a level uniformly on the whole lattice in
$\mathcal I$, and then average the sparse failure fraction over its
phase.  Up to the endpoint levels, this average is exactly the
typical-level failure probability.  Markov's inequality gives an
$o(N/L)$ number of unsuccessful probes.

We verify the extra separation, since it is needed when several
auxiliary charges occur on the same path.  The unscaled base choices
are independent of the sampled log level.  In a typical-scale limit,
the normalized separation of two distinct bases is therefore either
zero or infinite.  In the infinite case, requests in a common positive
$J$ are both viewed from its past, while requests in a negative or
terminal $J$ are both viewed toward its future.  The clearance through
far times for the later base in the former case, and for the earlier
base in the latter, includes the other base's time.  At that time an
actual list threshold is the other base's label.  Thus the two centers
are infinitely separated in the lifted interval.  Aperture-end
clearances also exclude proximity through the cyclic seam.  The fixed
ring order and its local linear comparison transfer this separation
to the circumference.  Decreasing the long margins by a fixed factor
gives the claimed finite separation.

All required factors can be imposed in one slow diagonal.  For example,
the unit-window oscillations in a near-link are finite in each bounded
test; $e^{\sqrt L}$ can be chosen as their increasing cutoff while the
larger separation scale is $e^{2L}$.  Fixed spatial and slot-rounding
enlargements can be included at the same time.  Omitting finitely many
additional levels at the probe ensures that every fixed factor
$CP(i)$ in a record requirement is at most $P(j+1)$ on its retained
assigned levels.  Auxiliary raw paths can be augmented beyond their
stopping times for these necessary tests.

The first-consumption window can be bounded by an absolute constant in
boundary units, using the zero-distance case of the passage exclusion,
the bounded endpoint conventions, and long-arc transmission.  Here
consumption of an original occurrence means reaching it as active or
using it as a paired or deleted slot.  Merely passing an untouched
interval into a fresh detached child is not consumption.  Once the
occurrence has passed strictly into such a child, the same main
trajectory in its starting polygon cannot later reach it.  Therefore
the ladder depth on the reaching side is its depth at the start of
that auxiliary path.  The array of middle labels in one $J$ is
determined by the original base heights and the fixed ring convention.

Choose actual witness times and types before the independent time-grid
translation used below; dummy or unsuccessful probes can be omitted.
Only finitely many tags per base and probe remain.  A choice of $J$
is recovered by the level-tracing rule from the main history, so it
does not supply an additional freely chosen main time.  This gives
the stated tag count.
\end{proof}

\subsection{Address codes, large moves, and time clusters}

Translate all nested time grids, of widths $r(j)$, by one common
uniform translation in $[0,T]$.  A bounded number of root boxes covers
the timeline.  Encode each base down to level $J_0$, and each near-link
time $q$ down to its probe level $i$.  An address is active at level
$j$ if its encoding depth is at least $j$.  Let $m_j$ be the number
of boxes occupied by active addresses, and set
\begin{equation}\label{addr:volume-code}
                         V=g\sum_{j=0}^{J_0}m_j.
\end{equation}
In particular $V\ge cN$.  There are
$n_0=O_d(1+N/L)=o(N)$ labeled entries.  Bounds denoted by
$o(N+V)$ below are taken with $d,R$ fixed and are uniform in the
admissible total count $V$.

\begin{lemma}[Entropy of the address code]\label{addr:code-entropy}
Including the witness tags, the logarithm of the number of preliminary
codes with specified total count is at most
\begin{equation}\label{eq:original-20}
                       V+O(V/g)+o(N+V).
\end{equation}
The constant in the forest term is absolute.  There are only
polynomially many possible total counts.
\end{lemma}

\begin{proof}
The occupied boxes form an ordered forest with a bounded number of
roots.  At fixed node count its unlabelled shape has at most
exponentially many possibilities, by parenthesis coding.  Each child
has at most $R$ possible positions in its parent, contributing at
most $g$ per child.  These costs give $V+O(V/g)$.

Attach each row label to its terminal box at its probe depth, and
each base label to its box at the deepest level.  On the proposed
bad-time grid, a base needs only $O_R(1)$ further choices within that
box.  Only a bounded number of entries terminate at each level.  The
remaining logarithmic assignment cost is bounded by a sum of
$O_d(n_0)$ terms of the form $\log(Cm_j)$, with bounded
multiplicity.  Concavity bounds this sum by
$O_d(n_0\log(1+C V/n_0))=o(N+V)$.  The tags from
Lemma~\ref{addr:finite-tests} have the same small-order cost.  The
available number of nodes is polynomial in $N$ for fixed $d$, so
summing over its possible integer totals adds only a logarithmic term.
\end{proof}

We now estimate probabilities under the original law with a code held
fixed as a specification.  We are not sampling a walk conditioned to
satisfy its code.  The bounded balanced-segment likelihood from
Section~\ref{sec:endpoints} permits replacement by the raw iid walk
of Equation~\eqref{eq:original-12}.  If needed, extend that raw walk
past the original endpoints to supply complete testing slots.

For a near-link row $e$ at probe $i$, choose its representative
$\bar q$ in its terminal box as follows.  It is either the first
state index of that box lying in the timeline, or the state immediately
after a specified move in the box whose absolute size is greater than
$P(i)$.  For an actual witness, take the state after the last such
move before $q$, if there is one, and otherwise take the first state.
Let $G_e$ be the largest oscillation, in either coordinate, of partial
sums in the box after erasing all moves larger than $P(i)$ in
absolute value, divided by $P(i)$.  Set
\begin{equation}\label{addr:row-tolerance}
 D_e=\left\lceil\log(1+G_e)\right\rceil,
 \qquad w_e=CP(i)\bigl(e^{\sqrt L}+e^{D_e}\bigr).
\end{equation}
This tolerance suffices for the equality at the representatives and
for the outward-wall tests on intervening subslots away from the
terminal boxes.  In the count, $D_e$ is a deterministic bin code.

Group specified representative moves according to which are the same
move.  Call them forced moves.  Sum over distinct groups by the iid
factorial-moment bound: choose ordered distinct time indices and use
their ordinary single-move probabilities.  All other increments remain
iid in this upper bound.

\begin{lemma}[Erasure bins and forced moves]\label{addr:erasure-bins}
For any fixed small $\eps>0$, it suffices, up to negligible total
probability times address-counting mass, to retain codes satisfying
\begin{equation}\label{eq:original-21}
                             \sum_eD_e\le\eps(N+V).
\end{equation}
The bin choices and sharing partitions have logarithmic cost
$o(N+V)$.  A forced group whose finest assigned depth is $i_0$
and whose coarsest assigned depth is $i_1$ has integrated mass at
most $C\exp(-cg(i_0-i_1))$ for a fixed $c>0$.
\end{lemma}

\begin{proof}
Partition the specified terminal-box step intervals at all their
integer endpoints.  For a resulting interval of length $w\ge1$,
truncate ordinary moves at
$u_w=c_0w\log(w+2)$, with $c_0>0$ small enough that
$u_w\le P(i)$ for every box covering that interval.  The bounded
increments, truncated variance and drift, and tail bound in
Equation~\eqref{eq:original-12} imply bounded small exponential
moments for their maximal partial sum divided by $u_w$.  The number
of ordinary moves larger than $u_w$ has a bounded small exponential
moment as well.

In summing the $G_e$, the coefficients of the first normalized bound
over all boxes covering a fixed interval sum to $O_d(1)$.
Indeed the assigned scales are geometrically separated with bounded
multiplicity.  A larger move of size $x$ contributes at most
$O_d(1)$, since the relevant sum of coefficients is
$\sum_{P(i)\ge |x|}|x|/P(i)$.  A forced move contributes the same
bound individually.  Independence on disjoint ordinary intervals
therefore gives a small exponential moment of $\sum_eG_e$ whose
logarithm is $O_d(n_0)$, uniformly after fixing the forced choices.

There are $o(N)$ rows.  By concavity, violation of
Equation~\eqref{eq:original-21} forces
$\sum_eG_e=\omega(N+V)$, uniformly for fixed $\eps$.
The exponential bound thus costs $\exp(-\omega(N+V))$, before
any of the finer tests is imposed.  It absorbs arbitrarily large
fixed multiples of the code entropy in
Equation~\eqref{eq:original-20}.  Conversely, composition counting
for nonnegative integer bins with the retained total bound gives
$\exp(o(N+V))$ choices.

For a forced group, its time has $O(r(i_0))$ choices and its size
must exceed $P(i_1)$.  The logarithmic Cauchy tail in
Equation~\eqref{eq:original-12}, together with comparability of the
logarithms in our scale range, gives
$C\exp(-cg(i_0-i_1))$ after summing both time and size.
Sharing partitions do not require a permutation of all row labels.
Choose the first label of each group in a fixed order, then its
highest probe and a bounded-size mask of probe labels at each sparse
level across its span.  This uses
$O_d(1+\text{span in logarithmic units}/L)$ digits per group.
Part of the exponential span payment sums these masks and the
possible spans, leaving a total factor $\exp(O_d(n_0))$.
Distinctness constraints may be dropped when multiplying nonnegative
upper bounds.  This also preserves the uniform erasure-tail estimate
conditional on the forced choices.

The representative is used only as a witness for the equality.  We
never condition on the unknown small-step state inside its box that
originally furnished the link, and we pay no sum over those states.
The same argument discards events exponentially small in a positive
power of the smallest time scale, which will be used for allocation
imbalances below.
\end{proof}

Retain the full set of leaf positions: all representative times and
the exact base times, including entries that are inactive at a finer
level.  At level $j$, split this ordered set at gaps greater than
$8r(j)$.  A time node is a resulting block containing an active
entry; let $n_j\le m_j$ be the number of nodes.  A node is regular
if it has an active center $s$ with no other full leaf at distance
in $[r(j)R^{-4},40r(j)]$.  The whole block is then within
$r(j)R^{-4}$ of $s$, and all other leaves are beyond its testing
slots.

\begin{lemma}[Regular time nodes]\label{addr:regularization}
For every fixed full set of leaves, the total number over all levels
of irregular occupied boxes is $O_R(n_0)$.  With arbitrarily high
probability under the lift over the extra grid translation, by taking
$R$ large and then the sequence parameters large, the witnessed
representatives also satisfy
\begin{equation}\label{eq:original-22}
 g\sum_j(m_j-n_j)\le O_g(n_0)+O(R^{-1})V.
\end{equation}
This restriction can be imposed in the forced-index sums; it need
not hold for all possible forced positions before they are specified.
\end{lemma}

\begin{proof}
At a fixed level, choose a sufficiently separated subcollection
containing a fixed fraction of the irregular boxes.  Each selected
box exhibits a pair of full leaves at a distance between the two
regularity cutoffs.  Starting with a maximum-cardinality packing of
the leaves at separation $200r(j)$, pass to separation
$cR^{-4}r(j)$.  The selection can be made sparse enough that, for
each pair, at least one new point may be added while retaining the
old packing.  Thus the number of irregular boxes is bounded by a
constant times the difference of these two packing counts.  Sum over
levels.  The two separation scales differ by a fixed number of
$R$-adic levels, so the differences telescope and give $O_R(n_0)$.
This also bounds the number of nonregular blocks.

First apply the random translation to the true witness times before
replacement by $\bar q$.  A regular tight cluster crosses a grid
division with conditional probability $O(R^{-4})$ and then occupies
at most two boxes.  The clusters are determined before the
translation.  Their active counts are bounded by $m_j$, and the
irregular contribution has the deterministic bound just proved.
Markov's inequality gives Equation~\eqref{eq:original-22} for these
true positions, with the weaker $O(R^{-1})$ loss.

Now apply the packing bound again to the complete representative set.
Consider a regular representative block split over two boxes.  Unless
the level is within a bounded distance of the encoding terminus of
a replaced entry, the corresponding true positions remain close and
in those same boxes.  They therefore belong to one true full-gap
block.  The representative split is already paid by the true
box-minus-block excess, because a box cannot contain anchors of two
distinct full-gap blocks.  Exceptional levels near encoding termini
cost $O_g(n_0)$.  This proves the displayed bound for the witnessed
representatives, without asserting it for unrelated forced choices.
\end{proof}

Put a graph on the $n_j$ time nodes, with active equality rows as
edges, retained through their encoding depth $i$.  A loop at one
node does not connect nodes.  Every component contains a base.
Let $u_j$ be the number of singleton components, $v_j=n_j-u_j$
the number of nodes in nonsingleton components, and $c_j$ the number
of nonsingleton components.  Let $l_j$ be the number of active rows
that are imprecise at this level, meaning $w_e>P(j+1)$.

\begin{lemma}[Tests selected by the code]\label{addr:necessary-tests}
The retained erasure codes satisfy
\begin{equation}\label{eq:original-23}
                   g\sum_jl_j\le C\eps(N+V)+o(N).
\end{equation}
At each regular node incident to a precise interblock edge, one may
require a directed survival bit toward its other endpoint, with
tolerance $CP(j+1)$ at the representative.  At each regular
singleton covered by a successful probe assignment, one may require
two bits, a singularity, or one main bit with its auxiliary charge.
Writing $u_j^{\rm cov}$ and $u_j^{\rm uncov}$ for the covered
and uncovered singleton counts, and treating irregularity losses
separately, one has
\begin{equation}\label{eq:original-24}
\begin{split}
 g\sum_ju_j^{\rm uncov}
 &\le CN+d\bigl[(a_+-1)+(b-b_-)\bigr]N+o_d(N),\\
 g\sum_ju_j^{\rm cov}&\le d(b-1)N+o_d(N).
\end{split}
\end{equation}
The constant $C$ in the first line may be chosen independently of
$d$ and $R$.
\end{lemma}

\begin{proof}
For a row attached at depth $i$, imprecision can persist only over a
logarithmic length bounded by $C D_e+O(\sqrt L+1)$, in view of
Equation~\eqref{addr:row-tolerance} and the scale ratios.  Sum over
rows and use Equation~\eqref{eq:original-21}.  Since there are
$O_d(N/L)$ rows, their $O(\sqrt L+1)$ contribution is $o(N)$.
This proves Equation~\eqref{eq:original-23}.

For a precise interblock row, the true interval has the required
outward wall, and replacement of its endpoint by its representative
changes time by less than $r(j+1)$.  The shell slot between, for
example, $4r(j+1)$ and $r(j)$ lies strictly between its two true
endpoints.  In fact precision and the $e^{\sqrt L}$ part of the
tolerance force the terminal-box length to be $o(r(j+1))$.
The regular center is within $r(j)R^{-4}$ of the representative.
Thus the wall test on that slot has at most a fixed multiple of
$P(j+1)$ additional allowance.  This reasoning also applies when
the representative lies immediately before or after a large move.

A covered singleton cannot have a near-link witness: its far
representative would lie outside the tight block and would connect
it to another node.  The other alternatives of
Lemma~\ref{addr:finite-tests} give precisely the stated tests.
Shifting from its base to the chosen regular center is absorbed in
the inner core.

For $gj\le N-O(1)$, all bases fit in at most one full-gap block.
This contributes $CN$ to the uncovered count.  Above that range
there are at most $d$ singleton components, because each contains a
base.  The unprobed portions have logarithmic lengths
$(a_+-1)N$ and $(b-b_-)N$.  Sparse-edge omissions, the common phase,
and unsuccessful probes remove only $o(N)$ additional log length.
This proves both bounds in Equation~\eqref{eq:original-24}.

Finally, fix ordered tie-breaking conventions for the regular center,
the row charged at an incident node, and the qualifying base at a
singleton.  These choices are determined once the tags, bins, and
forced positions are specified.  Codes that do not admit the
necessary tests contribute zero.  No arbitrary new bit assignment
at every occupied node and level is included in the count.
\end{proof}

\subsection{Auxiliary record charges}

Fix parameters
\begin{equation}\label{addr:moment-exponents}
                 0<\lambda<\frac12,
                 \qquad 2\lambda<\theta<\beta<1.
\end{equation}
The parameter $\lambda$ will eventually be chosen arbitrarily close
to $1/2$.  The strict inequalities in
Equation~\eqref{addr:moment-exponents} are used for two simultaneous
survival bits, for oscillation moments, and for summation over noise
bins, respectively.

The next three subsections establish the probability cost of a fixed
address code. The target negative logarithmic gain is
\[
 g\sum_j\left[(v_j-c_j)+\lambda v_j
                    +\frac32\lambda u_j^{\rm cov}\right],
\]
up to the errors from irregular nodes, imprecise rows, and sparse omissions.
The term $v_j-c_j$ comes from independent equality constraints between time
nodes. Each nonsingleton node also supplies a survival bit, costing
$\lambda$. A covered singleton supplies either two main bits, a singularity,
or a main bit plus an auxiliary test. The last alternative is the weakest:
its auxiliary test supplies the additional $\lambda/2$. We first prove
that auxiliary estimate, then extract the equality gain and integrate it
jointly with the main shell tests.

\begin{lemma}[A descending-record bound]\label{addr:record-bound}
For one coordinate of the thinned raw signed-move walk, the probability
that a descending record depth enters $[x-w,x+w]$ is at most
$C_\lambda(w/x)^\lambda$, uniformly when $w\gg1$ and $x\gg w$.
For several ordered depths, write $d_i'$ for the gap from the preceding
depth, taking the first gap from zero.  If the widths $w_i$ of both
neighbors are small compared with each gap, the probability of all
the required record events is at most
\[
                       \prod_i C_\lambda(w_i/d_i')^\lambda.
\]
\end{lemma}

\begin{proof}
Set $F(s)=(1+s)^{-\lambda}$ for $s\ge0$ and $F(s)=0$ for
$s<0$.  This is eventually superharmonic for the signed walk.  To
verify this, divide its one-step expected increment by $F(s)$ and
accelerate time by order $s/\log s$.  The density equivalents in
Equation~\eqref{eq:original-12} give, up to a positive constant,
\begin{equation}\label{addr:superharmonic-integral}
\begin{split}
 \PV\int_{\R}
 \frac{(1+t)^{-\lambda}\ind_{\{1+t>0\}}-1}{t^2}\dd t
 &=-1+\int_0^1
       \frac{z^{-\lambda}+z^\lambda-2}{(1-z)^2}\dd z\,<0.
\end{split}
\end{equation}
At $\lambda=1/2$ the last integral equals one, as follows by the
substitution $z=u^2$.  Its integrand is strictly increasing with
$\lambda>0$, so the strict sign follows for $\lambda<1/2$.
The passage from the walk to the integral is justified at small
increments by the bounded signed truncated drift and a Taylor
remainder.  Near the landing singularity and at an exact landing at
zero, the point-mass bound $C\log j/j^2$ gives the required
integrable domination.  Away from these regions the density
equivalents and the tail bound apply directly.

Start at height $x+w$ and stop on first descent to height at most
$2w$.  Before stopping, the heights are large enough for
superharmonicity.  Nonnegative optional sampling shows that the
probability of landing nonnegatively is at most
$C F(x+w)/F(2w)\le C_\lambda(w/x)^\lambda$.
A record in the requested depth interval requires such a landing.
For several separated target depths, restart at each successive
first descent.  The starting distance for the next target is
comparable to its gap $d_i'$, because the neighboring widths are
small compared with that gap.  Multiplication of these conditional
bounds proves the second assertion.

For application to both sides, generate two independent infinite
signed-move sequences and interleave them using independent fair side
choices.  Ignoring the segment stopping rule for the necessary record
events leaves the two thinned sequences independent.  This does not
assert independence of later placement histories from the segment;
those histories will be handled conditionally.
\end{proof}

\begin{lemma}[Conditional auxiliary gain]\label{addr:auxiliary-charge}
For the prescribed auxiliary tests at qualifying regular singletons,
the conditional probability cost is at least
$\lambda g/2-O(1)$ per charged level, in negative logarithmic
units, with an additional total factor $\exp(O_{d,R}(n_0))$.
The estimate holds after conditioning on the whole main word and
its seam parameters, and can be multiplied in the prescribed order
of the auxiliary explorations.
\end{lemma}

\begin{proof}
Condition on the main word, its seams, the code, and the representative
positions, without inspecting filling interiors.  The tags determine
at most $O_d(n_0)$ relevant paths, including predecessors needed to
place them. These are fixed indices in the prescribed seed schedule: an
unavailable start is skipped without renumbering any later seed. A tag
requiring that skipped path contributes zero. Group requests in each filling $J$. At the start of
a path, its original-label positions, long-arc availability, and
directed depths are known from earlier operations.  If those
conditions do not hold, the event contributes zero.

On one side of a path there are only $O(d)$ distinct label
prescriptions, regardless of how many scales charge them.  At a
fixed level $j$, distinct regular singleton requests are separated
in directed depth by $\gg_{d,R}P(j)$ and are that far from zero.
Indeed they have a common probe scale and the long mutual arc
separations of Lemma~\ref{addr:finite-tests}, with the arcs still
pending in the starting polygon.  An origin is charged at most once
at that level.  The indicated constants need only exceed sufficiently
large constants depending on $d,R$.

For each obstruction depth retain its smallest required width
$w=CP(i)$.  Process these depths in increasing width, deleting a
candidate if it is within $M_1$ times its width of an already kept
depth, where $M_1$ is fixed and large compared with $R$.
Then order the retained depths spatially.  Their gaps satisfy the
width separation required in Lemma~\ref{addr:record-bound}.
At level $j$, each requested charge is assigned to a retained depth
whose width is at most $P(j+1)$ and whose distance from the requested
depth is at most $M_1P(j+1)$; the omitted levels near a probe ensure
this assertion.  If its preceding gap is less than $P(j)$, move to
its predecessor and repeat.  Every representative encountered still
has width at most $P(j+1)$.  The separation between charges and from
zero, and the bounded number of prescriptions per side, ensure that
distinct arriving charges terminate at distinct retained depths with
preceding gap at least $P(j)$.

Consequently the sum $\sum_i\log(d_i'/w_i)$ pays at least $g$
per obstruction charge, summed over all levels, up to harmless
rounding constants per path.  Lemma~\ref{addr:record-bound} gives
the corresponding obstruction gain $\lambda g$.
For origin bits or singularities, the one-sided version of the shell
estimate and nested integration proved below in
Lemmas~\ref{addr:shell-moment} and~\ref{addr:joint-integration}
gives $\lambda g-O(1)$ per level.  That special case has only a
forward window about time zero, an augmented iid path, and no forced
moves or smoothing chunks.  Its proof is independent of the present
auxiliary estimate.  The balanced trajectory likelihood costs a
bounded factor per path.

For a path bearing both kinds of tests, its joint probability is at
most the smaller of their separate upper bounds and hence at most
their geometric mean.  It therefore pays
$\lambda g/2-O(1)$ per auxiliary charge of either type.  The
$O(1)$ per origin-bit shell may be chosen bounded as $R$ grows.
Constants per prescription and per path give the stated
$\exp(O_{d,R}(n_0))$ factor.

Finally integrate the fresh polygon interiors backwards in their
prescribed exploration order.  Keep the indicators of later
reachability and placement constraints until their starting histories
are known.  Each bound is uniform over the applicable starting
histories, and the number and tags of the charges on each grouped
path have already been fixed by the main data and representative
positions.  Thus the conditional estimates multiply without assuming
independence of those histories.
\end{proof}

\subsection{Equality smoothing and its rank gain}

In every gap of the full ordered time list of length
$s\ge r(J_0)/2$, choose a noise chunk in a middle portion, of
length comparable to $s$.  These chunks are disjoint.  A forced
move is within distance one of a full representative, so no chunk
contains a forced move.  Initially condition on all nonchunk
increments.  Each chunk contributes two displacement columns.
For either column $\ell$ let
\[
                  a_\ell=\frac{2}{\pi^2}s\log s,
\]
and bin its coordinate's maximal oscillation divided by $a_\ell$
in upper dyadic bins $B_\ell\ge1$.

\begin{lemma}[A point-mass bound with a range bin]\label{addr:chunk-density}
Conditional on typical binomial allocation counts between the two
sides in each chunk, the joint lattice point mass of its displacement
columns in the specified range bins is at most
\begin{equation}\label{addr:chunk-mass}
                  \prod_\ell C_\beta a_\ell^{-1}
                                      B_\ell^{-1-\beta}.
\end{equation}
Exponentially atypical side-count imbalances have negligible total
mass in the address estimate.
\end{lemma}

\begin{proof}
For one coordinate, the stable lattice local limit estimate on a
comparable number of signed moves gives the uniform point-mass
bound $C/a_\ell$.  Let $B=B_\ell$ be large and fix a sufficiently
small $\zeta>0$.  If a move exceeds $B^{1-\zeta}a_\ell$, then,
conditioning on all other moves and summing over its possible time,
its contribution to any specified displacement is at most
\[
 Cs\sup_{j>B^{1-\zeta}a_\ell}\frac{\log j}{j^2}.
\]
Because $a_\ell$ is comparable to $s\log s$, this is bounded
by $C_\beta a_\ell^{-1}B^{-1-\beta}$ after choosing
$\zeta$ small relative to $1-\beta$.  Logarithmic losses are
absorbed by the strict power margin.

If there is no such large move, an oscillation of order $Ba_\ell$
requires an oscillation of that order on at least one half of the
chunk.  Maximal Bernstein, the truncated second moment, and the
truncated drift bound make that event smaller than any fixed inverse
power of $B$.  Apply the local limit estimate to the independent
other half to retain a factor $C/a_\ell$.  This gives the desired
point-mass and bin bound in the complementary case.

Given the allocation counts, or the allocation locations, the two
thinned signed sums and their ranges separate.  Different chunks
also use independent increments, proving the product bound.
Atypical allocations are exponentially unlikely in the chunk length.
The smallest chunk length is a positive power of the large parameter,
so Lemma~\ref{addr:erasure-bins} discards their total contribution
before the finer estimates.
\end{proof}

Each row compares a difference of two coordinate values at full
leaves.  It is therefore an integer linear function of the displacement
columns, plus a shift determined by nonchunk increments and forced
moves.  Write $M$ for this matrix; its entries are $0,1,-1$.
A nonzero minor with selected row set $I$ and column set $J$ has
\emph{profit}
\[
                   \sum_{\ell\in J}\log a_\ell
                         -\sum_{e\in I}\log w_e.
\]
The empty minor is allowed and has profit zero.

\begin{lemma}[Threshold rank and minor profit]\label{addr:weighted-rank}
For any such finite matrix and positive widths and scales, the
largest profit of a nonzero minor is at least
\begin{equation}\label{addr:rank-integral}
 H_M=\int_{\R}
       \rank M[\log w_e\le t,\ \log a_\ell\ge t]\dd t.
\end{equation}
The selected minor can be taken in the original integer matrix $M$.
\end{lemma}

\begin{proof}
Put $u_e=\log w_e$ and $v_\ell=\log a_\ell$.  Order rows by
increasing $u_e$ and columns by decreasing $v_\ell$.  Elimination
using only earlier rows added to later rows, and earlier columns
added to later columns, reduces $M$ to a partial permutation matrix
up to nonzero pivot factors.  Explicitly, pivot on the earliest
nonzero entry in the first nonzero row, clear its later row and
column entries, and continue on the remaining ordered lists.
These operations preserve the ranks of all northwest prefix
submatrices.  A pivot at $(e,\ell)$ contributes one to the
integrand exactly on $u_e\le t\le v_\ell$.  Hence $H_M$
is the sum of $(v_\ell-u_e)_+$ over the pivots.

To pass back to a minor of the original matrix, introduce a parameter
$A\to\infty$, and multiply row $e$ by $e^{-Au_e}$ and
column $\ell$ by $e^{Av_\ell}$.  The triangular multipliers
in the elimination stay bounded after this conjugation: each
earlier-to-later row operation receives a factor
$e^{A(u_{\rm earlier}-u_{\rm later})}$, and each corresponding
column operation receives
$e^{A(v_{\rm later}-v_{\rm earlier})}$.  Both exponents are
nonpositive.  Matrix dimensions and coefficients are fixed before
$A$ increases.

The minor of the reduced scaled matrix consisting of the
positive-profit pivots has magnitude a nonzero constant times
$e^{AH_M}$.  Cauchy--Binet applied to the triangular products
bounds it by a constant times the largest scaled nonzero minor of
the original matrix.  Taking logarithms, dividing by $A$, and
letting $A\to\infty$ gives the claim.  If there are no
positive-profit pivots, the empty minor suffices.  No integrality
of the elimination multipliers is asserted or needed; integrality
will be used only for the selected original minor.
\end{proof}

\begin{proposition}[Equality gain]\label{addr:equality-smoothing}
After extracting equality constraints, the remaining bin weights
can be bounded by $\prod_\ell B_\ell^{-\beta}$, with a
factor $\exp(O(n_0))$.  The equality profit is at least
\begin{equation}\label{eq:original-25}
                           g\sum_j(v_j-c_j-l_j),
\end{equation}
with a nonnegative rank contribution used at levels where a summand
is negative.
\end{proposition}

\begin{proof}
Choose rows and an equal number of columns forming a nonzero minor
of the original integer matrix.  Fixing unselected columns gives
at most $\prod_{\ell\notin J}CB_\ell a_\ell$ possible
integer values in their bins.  For selected columns, the row slabs
contain at most $\prod_{e\in I}Cw_e$ integer points.  To see
this, enlarge each integer point to its unit cube.  The original
matrix has entries $0,\pm1$, so each row width grows by at most
$O(n_0)$, which is negligible compared with $w_e$ in our scale
range.  Change variables using the selected minor.  Its determinant
is a nonzero integer, of absolute value at least one.  The resulting
slab volume is bounded by the product of the enlarged widths.

Multiply these counts by Equation~\eqref{addr:chunk-mass}.
The selected columns and rows contribute the exponential of minus
their profit.  Each unselected column retains $B_\ell^{-\beta}$;
selected columns retain at least that decay as well.  The constant
cost is exponential in the number of rows and chunks, hence
$\exp(O(n_0))$.

Apply Lemma~\ref{addr:weighted-rank}.  For
$t\in[\log P(j+1),\log P(j)]$, all precise rows and all
gap columns separating the full time blocks are available.  Those
columns allow independent shifts in both coordinates of every
full-gap block, except possibly the first.  In the graph on active
nodes, a forest has $v_j-c_j$ independent connecting rows.
Independence holds even when edges use different coordinates:
remove a nonroot leaf and use the coordinate appearing in its
unique incident row.  If a block has no available shift, take it
as root.  Removing the $l_j$ imprecise rows loses at most $l_j$
units of rank.  The threshold rank is thus at least
$(v_j-c_j-l_j)_+$.  The logarithmic width of each scale band satisfies
\[
 \log P(j)-\log P(j+1)
   =g+\log\frac{\log r(j)}{\log r(j+1)}\ge g.
\]
The gaps greater than $8r(j)$ leave a fixed margin for the middle
chunks and their integer rounding.  Summing these rank lower bounds
over the bands proves Equation~\eqref{eq:original-25}.
\end{proof}

\subsection{Shell integration in the presence of rare jumps}

Around a tested regular center $s$ at level $j$, use the window
and inner core
\begin{equation}\label{addr:windows}
 I=[s-r(j),s+r(j)],\qquad
 I_*=[s-4r(j+1),s+4r(j+1)],
\end{equation}
with integer rounding on the timeline.  Regularity is tested against
all full leaves, not only the active ones.  It follows that distinct
test windows are disjoint unless the smaller lies in the inner core
of the larger.  Their shell slots contain neither forced moves nor
noise chunks.  Indeed gap endpoints relative to a regular center are
either very near or beyond the tested window.  A middle portion of a
gap extending away from its tight cluster is already beyond the shell;
one may choose middle fourths for the chunks.  Thus any chunk meeting
a tested window is wholly in its inner core.  For an auxiliary
origin test, use the corresponding forward window and core.

\begin{lemma}[Shell moment with survival]\label{addr:shell-moment}
Condition on the non-shell data and let $D_*$ be a fixed nonnegative
upper bound on the oscillation in the inner core under this conditioning.
Let $D$ be $D_*$ plus the sum of the maximal shell oscillations.
For $k'\in\{1,2\}$ distinct requested survival bits, iid sampling
of the shell increments gives
\begin{equation}\label{eq:original-26}
 \E\!\left[
    1+\bigl(D/P(j)\bigr)^\theta\,;\ 
       \text{the bits with their allowance}
       \right]
 \le C R^{-\lambda k'}
       \left(1+\bigl(D_*/P(j+1)\bigr)^\theta\right).
\end{equation}
For a singularity test the factor $R^{-\lambda k'}$ may be
replaced by $R^{-\theta}$, with the necessary restriction
$D_*\ge cP(j)$ on its input.  The constants can be bounded as
$R$ increases.  The limit order is to fix $R$ first and then
take the large sequence parameters.
\end{lemma}

\begin{proof}
The allowance at the inner endpoint of a shell bit is at most
$D_*+CP(j+1)$.  When $D_*/P(j+1)\gtrsim R$, ordinary maximal
$\theta$-moments suffice.  Indeed the right side of
Equation~\eqref{eq:original-26} dominates both the constant term
and $(D_*/P(j))^\theta$, since $\theta>2\lambda\ge\lambda k'$.

For the remaining values, fix $R$ and pass to the symmetric Cauchy
limits of the shell increments.  Their durations, in units of
$r(j)$, are bounded above and below by positive constants.  Distinct
coordinates and the two time directions separate in this limit.
For a symmetric Cauchy process $X$ and $0<u\le C$, the probability
of remaining above $-u$ for a unit time is at most
$C_\lambda u^\lambda$.  Here is a direct way to obtain the
exponent used.  At an independent exponential time of rate one, the
classical Wiener--Hopf identity
\citep[Chapter~VI, Theorem~5 and Corollary~10]{BertoinLevy} gives the logarithm
of the Laplace transform of the supremum as
\begin{equation}\label{addr:wh-transform}
 \int_0^\infty e^{-t}
   \E\!\left[(e^{-aX_t}-1)\ind_{\{X_t>0\}}\right]
                  \frac{\dd t}{t}.
\end{equation}
The Cauchy scaling $X_t\overset{\rm law}=tX_1$ bounds this
expression by $-\lambda\log a+C_\lambda$ for every
$\lambda<1/2$.  To see the coefficient, integrate over
$t$ between order $a^{-1}$ and one: the integrand in the
expectation approaches $-1/2$, while the small- and large-time
remainders are controlled after choosing any strict exponent
$\lambda<1/2$.  A Laplace bound with $a=u^{-1}$ gives the
exponential-time survival estimate.  Survival to deterministic
time one is bounded above by a constant times exponential-time
survival, by integrating over exponential times in $[0,1]$.
Symmetry changes the supremum statement to the stated lower-wall
statement.

We also need its oscillation moment on the same survival event.
For $z\ge1$, the probability of survival together with a unit-time
maximal oscillation exceeding $z$ is at most
$Cu^\lambda/z$.  Decompose time into
$[2^{-m-1},2^{-m}]$, $m\ge0$.  A large total oscillation requires
an oscillation of order $z2^{-\eta m}$ on some such interval,
where $\eta>0$ is small.  The symmetric maximal inequality gives
cost $O(2^{-m(1-\eta)}/z)$ there.  Survival before that interval
has probability at most $C(u2^m)^\lambda$ and depends only on
earlier increments.  Multiplication and summation are allowed
because $1-\eta-\lambda>0$.  Other independent bits contribute
their own survival factors.  Integrating the resulting $1/z$ tail
gives the $\theta$-moment bound, since $\theta<1$.  An untested
shell coordinate can be integrated independently.

Apply these estimates with
$u=C(1+D_*/P(j+1))/R$.  If the inner bound is fixed, each
survival process starts at its own shell's inner endpoint, with
the left process reversed when required.  The wall measured from
the leaf implies survival for these shell increments with the
stated allowance.  The moment of the shell maxima under all bits
is at most a constant times $u^{\lambda k'}$.  The transmitted
term $(D_*/P(j))^\theta$ obeys the same desired upper bound.
Because $\theta>\lambda k'$, these observations give
Equation~\eqref{eq:original-26} for the limiting processes.

For the walks, enlarge positive allowances by a fixed factor before
using convergence.  Their maximal $\theta$-powers are uniformly
integrable by the tail, truncated variance, and truncated drift
bounds in Equation~\eqref{eq:original-12}.  At fixed $R$, any
sequence of bounded input allowances has a convergent subsequence,
so the upper estimate is uniform over those allowances.  The
unbounded-input case was already handled by ordinary moments.
This proves the estimate with constants bounded as $R$ increases,
followed by sequence parameters sufficiently large depending on
$R$.  Finally a singularity requires
$D_*\ge cP(j)$; ordinary shell moments then give the same bound
with factor $R^{-\theta}$.
\end{proof}

\begin{lemma}[Joint integration of chunks, shells, and forced moves]
\label{addr:joint-integration}
The equality gain in Equation~\eqref{eq:original-25} can be
extracted simultaneously with the shell gains in
Equation~\eqref{eq:original-26}.  After summation over chunk bins,
forced moves, and their sharing partitions, the additional
logarithmic cost is
$O_{\lambda,\theta}(V/g)+o_d(N+V)$.
The statement also holds for forward origin windows without chunks
or forced moves, without using any auxiliary-charge estimate.
\end{lemma}

\begin{proof}
Use an upper bound for each inner-core oscillation that depends on
chunk increments only through their bins.  Namely, sum the bounds
$D$ from directly enclosed tested windows, the binned chunk maxima
outside those windows, the absolute forced sizes there, and the
maximal range of the concatenation of all remaining ordinary inner
increments outside these holes.  This sum dominates the true
oscillation.  The last concatenated range has one ordinary bounded
normalized moment at scale $P(j+1)$, because its total step
count is at most a constant times the inner-core length.

Apply Equation~\eqref{eq:original-26} successively to the nested
windows.  It passes a child weight
$1+(D_{\rm child}/P(j_{\rm child}))^\theta$ to the next
integration, charges a factor $1+CB_\ell^\theta$ for each chunk
at its first enclosing core, and charges
$1+C(|x|/P(j+1))^\theta$ for each forced move at its tightest
enclosing core, if there is one.  A chunk's length is bounded by a
constant times that core's length at the first enclosing fit.
Since $\theta<1$, the fractional moment of a sum is bounded by
the sum of the fractional moments; passing between this sum and
the corresponding nonnegative product costs only constant factors
per child, window, and nuisance contribution.  The total extra
shell constants therefore have logarithm
$O_{\lambda,\theta}(V/g)$.

For clarity, equality smoothing is performed in the same
calculation, not on a separately conditioned shell law.  First fix
codes, forced choices, and bins.  Replace the original tests by
the necessary enlarged tests just described, which see chunks only
through the bins.  Conditional on nonchunk increments, extract the
integer equality profit and the point-mass/bin weights of
Proposition~\ref{addr:equality-smoothing}.  The remaining
nonchunk increments still have their original iid integration.
Integrate the outer shell inequalities, which leave the displayed
child weights, and proceed inwards.  Every dyadic bin is then
summable, because
$B_\ell^{-\beta}(1+CB_\ell^\theta)$ is summable over upper
dyadic bins when $\theta<\beta$.  The forward origin case is
exactly the same nested integration with no equality holes or
forced moves and hence does not invoke
Lemma~\ref{addr:auxiliary-charge}.

It remains to sum forced moves with the propagated moments.
For a forced group with finest and coarsest depths $i_0,i_1$,
logarithmic comparability over the tested range bounds its propagated
factor by
\[
 C_{g,\delta}\left(1+\frac{|x|}{P(i_0)}\right)^\theta
 \left[
 \max_{\substack{j>i_0:\\
        \text{its time lies in a }Cr(j)\text{ active-box neighborhood}}}
        \left(\frac{r(i_0)}{r(j)}\vee1\right)
 \right]^\theta.
\]
The empty maximum is interpreted as one.  The size moment over
$|x|>P(i_1)$ still has an exponential span payment: before the
time sum it is bounded by
\[
                 r(i_0)^{-1}C e^{-cg(i_0-i_1)}
\]
for some $c>0$.  This follows from the logarithmic Cauchy tail
and $\theta<1$; the strict moment inequality leaves a positive
span exponent.

In the normalized time-position sum over the finest assigned box,
the expectation of the second factor is at most
\begin{equation}\label{addr:forced-position}
                  C\left(
                    1+\sum_{j>i_0}
                       R^{-(1-\theta)(j-i_0)}m_j
                    \right).
\end{equation}
Indeed the union of enlarged active-box neighborhoods at level
$j$ has at most $Cr(j)m_j$ integer positions.  This count is
uniform in all other exact representative choices.  Multiplying
its relative size by $(r(i_0)/r(j))^\theta$ gives precisely
the geometric term in Equation~\eqref{addr:forced-position}.

The majorants used in the different forced-position sums depend on separate
time variables. To see this explicitly, for the fixed preliminary code let
$U_j$ be the union of the $Cr(j)$-neighborhoods of its occupied active
level-$j$ boxes. These sets and $m_j$ are fixed before any exact forced
position is summed. Every tested center belongs to one of those boxes, so
a forced time $q$ in its inner core belongs to $U_j$. The position factor
for a group of finest depth $i_0$ is consequently bounded by
\[
 F_{i_0}(q)=\max\left\{1,\sup_{j>i_0}
               R^{\theta(j-i_0)}\ind_{U_j}(q)\right\}.
\]
This depends only on that group's time and the fixed code. Dropping
distinctness and other admissibility restrictions in the nonnegative upper
bound makes the sums of products factor over groups. The bound
$|U_j|\le Cr(j)m_j$ then gives Equation~\eqref{addr:forced-position}
for each factor.

There are $O_d(n_0)$ groups and bounded multiplicity of their
$i_0$ values at each level.  Sum the geometric convolution over
these values, then use concavity for the logarithm of its product.
Its logarithm is bounded by
$O_d(n_0\log(1+C V/n_0))=o(N+V)$.
Use the remaining exponential span payment to sum sharing
partitions exactly as in Lemma~\ref{addr:erasure-bins}.
These integrations use deterministic total exponents whenever
Equations~\eqref{eq:original-22} and~\eqref{eq:original-23},
the test coverage, and the stated separations hold.  They do not
require the exact cluster shapes to be independent of forced
positions.  This completes the simultaneous estimate.
\end{proof}

\subsection{The address contradiction}

\begin{proposition}[Completion of the bad-grid estimate]
\label{addr:bad-grid-proof}
The bad-grid bound in Proposition~\ref{end:bad-grid} holds, with
every prescribed failure exponent and in all the parameter ranges
specified there.
\end{proposition}

\begin{proof}
Suppose it fails along a large-parameter sequence.  As in
Section~\ref{sec:endpoints}, choose an interval with more than
$\exp(\tau(b-1)N)$ bad grid times on an event of probability
at least $e^{-B'N}$.  Condition on that event and sample $d$
bad times with replacement.  Equation~\eqref{eq:original-19}
shows that an event of lift probability bounded below has
counting-measure mass at least a fixed positive constant times
\begin{equation}\label{addr:counting-lower}
                      \exp\{d\tau(b-1)N-B'N\}.
\end{equation}
The constant $B'$ can absorb the number of intervals; it suffices
to obtain a contradiction for arbitrarily large fixed $B'$.

Impose the high-lift-probability requirements of
Lemma~\ref{addr:finite-tests} and the random grid restriction in
Equation~\eqref{eq:original-22}.  Discard the negligible bin and
allocation exceptions of Lemmas~\ref{addr:erasure-bins}
and~\ref{addr:chunk-density}.  We estimate the remaining mass
by summing its codes under the original law.

Main shells charge $\lambda g$ per directed bit and $\theta g$
per singularity.  On nonsingleton nodes these charges cover all
nodes except irregular ones and at most $2l_j$ endpoints of
imprecise edges.  On a qualifying regular singleton, two bits
or a singularity give at least $2\lambda g$, while a main bit
and its auxiliary test give at least
$\frac32\lambda g-O_{\lambda,\theta}(1)$ by
Lemma~\ref{addr:auxiliary-charge}.  Thus the latter is a common
lower gain for every covered singleton.

Combine these charges with equality smoothing and the joint
integration.  The negative logarithm of the integrated mass for
one code is at least
\begin{equation}\label{addr:charge}
\begin{split}
 g\sum_j\left[
       v_j-c_j+\lambda v_j+\frac32\lambda u_j^{\rm cov}
            \right]
 &-C(\eps+R^{-1})(N+V)\\
 &-O_{\lambda,\theta}(V/g)-o_d(N+V).
\end{split}
\end{equation}
Here irregular-node losses are $O_R(n_0)$, imprecise-edge losses
are controlled by Equation~\eqref{eq:original-23}, and
Equation~\eqref{eq:original-22} replaces node counts by box
counts at the displayed cost.  Constants multiplying $\eps$
can be absorbed by reducing its choice.  The auxiliary
separation and precision assumptions restrict only the applicable
histories for the specified singleton tests; no further choices
are introduced in summing the code.

Subtract Equation~\eqref{addr:charge} from the entropy in
Equation~\eqref{eq:original-20}.  At the limiting bookkeeping
value $\lambda=1/2$, the main expression is
\[
 g\sum_j\left[
        -\frac{v_j}{2}+c_j
        +\frac14u_j^{\rm cov}+u_j^{\rm uncov}
          \right].
\]
Each component contains a base, so $c_j+u_j\le d$.
Also a nonsingleton component has at least two nodes.  Hence
\[
             \frac{v_j}{2}-c_j
                    \ge \frac12(n_j-2d)_+.
\]
Since $gJ_0=bN+O(g)$, Equations~\eqref{eq:original-22}
and~\eqref{eq:original-24} yield the following upper bound
for the logarithm of the total retained mass at any admissible
count $V$:
\begin{equation}\label{addr:counting-upper}
\begin{split}
 &CN+d\bigl[(a_+-1)+(b-b_-)\bigr]N
       +\frac14d(b-1)N\\
 &\qquad-\frac12(V-2dbN)_+
       +\eps'(N+V)+o_d(N+V).
\end{split}
\end{equation}
There are polynomially many possible counts, so their summation
adds only a small-order term.  For $\eps'$ small, the negative
excess-$V$ term also controls arbitrarily large $V$; maximizing
over $V$ costs only the displayed errors at $V=O_d(N)$.

The error $\eps'>0$ can be made arbitrarily small by first
taking $\lambda$ close enough to $1/2$, choosing compatible
$\theta,\beta$, then taking $R$ large and the bin tolerance
small, with a sufficiently slow sparse stride.  The absolute
$C$ in the $CN$ term is independent of $d$.
Choose $a_+-1$ and $b-b_-$ small relative to
$(\tau-1/4)(b-1)$, then $d$ large relative to $B'+C$,
and finally all remaining errors sufficiently small.  The
coefficient of $N$ in Equation~\eqref{addr:counting-upper}
is then strictly smaller than
$d\tau(b-1)-B'$.  This contradicts
Equation~\eqref{addr:counting-lower} on the requirements having
lift probability bounded below, including the independent random
shifts.  Their probability can be made arbitrarily high by the
prescribed choices.  The bad-grid estimate follows.
\end{proof}

\subsection{The common joint law and deterministic calibration}

\begin{proposition}[Joint subsequential comparison]
\label{addr:subsequential-joint}
Put $b_n=\mathcal B(B_n)$.  From every sequence of positive
integers tending to infinity one can extract a further sequence
on which, in the prescribed flag uniformization,
\[
 \bigl((\phi_n)_*m_n,K_n(b_n^{-1}),\Gamma_n\bigr)
\]
converges jointly to
\[
 \left(
   \mu_h^{\rm crit},
   \{(z,w,cD_h(z,w)):z,w\in\widehatC\},
   \Gamma
 \right),
\]
where $0<c<\infty$ is deterministic.  The law of the marked
surface is the law constructed in Theorem~\ref{thm:continuum},
and $\Gamma$ is the independent nested $\CLE_4$ of
Theorem~\ref{thm:main}.  The maximal spherical flag mesh tends
to zero jointly on these extractions.
\end{proposition}

\begin{proof}
Proposition~\ref{addr:bad-grid-proof} supplies the estimate
required by Proposition~\ref{end:endpoint-reduction}.
The latter gives the all-endpoint modulus and macroscopic
diameter bounds.  Its temporary upper perimeter cutoff
$p_*=C_*B_n$ can be replaced in these statements by $B_n$:
for fixed $C_*$, eventual adjacent-scale bounds from
Proposition~\ref{metric:scale-bootstrap} keep
$\mathcal B(p_*)/\mathcal B(B_n)$ bounded.  The cutoff and
other truncations can then be exhausted in probability.

On a common profile extraction, ordinary open patches have
the same deterministic lower and upper factor by
Proposition~\ref{metric:scale-bootstrap}.  Endpoint control
extends their equality to every vertex pair as in
Proposition~\ref{end:endpoint-reduction}.  In particular a
bounded-cost sequence of paths cannot acquire an additional
limiting distance at an exceptional endpoint.  The factor is
strictly positive and finite, and is deterministic on the
chosen profile extraction.

All estimates first made under restricted central intensities
transfer to the exact finite sphere law by the volume local
limit domination and counting exhaustion in
Proposition~\ref{cut:fixed-volume} and
Corollary~\ref{cut:unbiased}.  The common conditioning kernels
are retained at fixed sphere volume by
Proposition~\ref{kern:fixed-volume}.  Thus the endpoint
comparisons join the area and loop comparisons on the same
probability space.  Proposition~\ref{conf:flag-uniformization}
identifies the actual discrete conformal uniformization and
gives vanishing flag mesh there.

We spell out the distance topology.  Every point of the sphere
lies in the image of a flag, and each flag contains an original
primal vertex.  Vanishing mesh therefore makes the embedded
primal vertices spatially dense.  For every limiting pair
$(z,w)$, choose approaching vertices.  The endpoint modulus
and the upper and lower comparisons identify their distance
as $cD_h(z,w)$.  Conversely, for any approaching sequence
of vertex pairs, the same modulus and comparisons rule out
an additional limiting distance value.  Diameter tightness
prevents escape in the third coordinate.  These are precisely
the two inclusions needed for Hausdorff convergence of the
full compact distance graph.

Area convergence is weak convergence in these coordinates.
For loops, the full-tour comparisons and count exhaustion
give partial bijections matching every loop above each
positive spherical diameter cutoff, with multiplicity one
and all descendants retained.  Their uniform individual-loop
errors tend to zero; unmatched loops are discarded only
below the corresponding vanishing cutoff.  This is the
loop topology of Theorem~\ref{thm:main}.  Finally, the
three-point normalization in
Proposition~\ref{conf:flag-uniformization} is the
area-marked law of Theorem~\ref{thm:continuum}, with
the independent nested CLE already identified by the
cutting and unbiased comparison.  All the assertions
therefore hold in one common joint law.
\end{proof}

\begin{lemma}[Bounded deterministic calibration]\label{addr:calibration}
There is a deterministic sequence $a_n>0$ tending to zero
such that the joint convergence asserted in
Theorem~\ref{thm:main} holds through all positive integers.
\end{lemma}

\begin{proof}
For each of the first two normalization samples, take the
original primal vertex of the flag containing it.  Samples
on flag edges have probability zero, so this is unambiguous
almost surely; any fixed convention handles that null set.
Let $Y_n$ be the original all-edge graph distance between
the two selected vertices.  Set
\[
                   Z=D_h(0,1),\qquad t_*=\E[\tanh Z].
\]
The continuum metric is finite and separates distinct
points, so $0<Z<\infty$ almost surely and $0<t_*<1$.

Along any further extraction of
Proposition~\ref{addr:subsequential-joint}, the two selected
vertex images approach $0$ and $1$, since each lies in
the flag containing its normalization sample.  Joint mesh
and distance-graph convergence imply
\begin{equation}\label{addr:sampled-distance}
                         Y_n/b_n\ \Longrightarrow\ cZ.
\end{equation}
In particular $\Prob(Y_n=0)\to0$ through the full sequence:
otherwise an offending subsequence would admit a further
extraction contradicting the positive limit in
Equation~\eqref{addr:sampled-distance}.

For all sufficiently large $n$, the function
$F_n(a)=\E[\tanh(aY_n)]$ is continuous and strictly increasing
for $a>0$, starts at zero, and tends as $a\to\infty$ to
$\Prob(Y_n>0)>t_*$.  Thus there is a unique deterministic
$a_n>0$ satisfying
\begin{equation}\label{addr:normalization}
                         \E[\tanh(a_nY_n)]=t_*.
\end{equation}
Assign arbitrary positive values to the finitely many initial
exceptions.  Coincident selected vertices at some finite
indices cause no problem for this argument.

On an extraction with factor $c$, bounded weak convergence
in Equation~\eqref{addr:sampled-distance} gives, for each
fixed $u>0$,
\[
           \E[\tanh(uY_n/b_n)]\longrightarrow\E[\tanh(ucZ)].
\]
The limiting function is strictly increasing.  Bracket
$1/c$ by any two fixed positive values, one below and one
above it, and use Equation~\eqref{addr:normalization}.
It follows that $a_nb_n\to c^{-1}$ on that extraction.
This uses no convergence of unbounded distance moments.

For any fixed permitted $\sigma>0$, Equation~\eqref{eq:original-18}
with $P'=1$ gives
\[
       \mathcal B(P)\ge
              \frac{\mathcal B(1)}{K_\sigma}P^\sigma.
\]
Since $B_n=(2/\pi)\sqrt n\log n\to\infty$, also $b_n\to\infty$.
Every sequence has a further extraction with
$a_nb_n\to c^{-1}\in(0,\infty)$; consequently $a_n\to0$
through the full sequence.  Multiplying the limiting
distance graph on that extraction by $a_nb_n$ removes $c$.
All resulting subsequential joint laws are the same law
specified in Theorem~\ref{thm:main}, so the convergence
holds through all positive integers.

The calibration changes only the deterministic graph-distance
unit.  The conformal structure remains that of the specified
equilateral flags, and the metric remains the all-edge
distance on original primal vertices.  Auxiliary sampled
topological positions have served only in the comparisons.
The area, complete nested loop collection, and vanishing
flag mesh remain in the same joint law throughout.
\end{proof}

Theorem~\ref{thm:continuum} and Lemma~\ref{addr:calibration}
complete the proof of Theorem~\ref{thm:main}.
